% Generated draft: requires mathematical and source review.
\documentclass[11pt]{article}
\usepackage[margin=25mm]{geometry}
\usepackage{fontspec}
\setmainfont{Latin Modern Roman}
\usepackage{amsmath,amssymb,mathtools,unicode-math}
\setmathfont{Latin Modern Math}
\usepackage{enumitem,needspace,xurl,hyperref}
\setlength{\parindent}{0pt}
\setlength{\parskip}{5pt}
\setlength{\emergencystretch}{3em}
\newcommand{\N}{\mathbb{N}}
\newcommand{\Z}{\mathbb{Z}}
\newcommand{\Q}{\mathbb{Q}}
\newcommand{\R}{\mathbb{R}}
\newcommand{\C}{\mathbb{C}}
\newcommand{\F}{\mathbb{F}}
\newcommand{\A}{\mathbb{A}}
\newcommand{\PP}{\mathbb{P}}
\newcommand{\E}{\mathbb{E}}
\newcommand{\eps}{\varepsilon}
\newcommand{\dd}{\,\mathrm{d}}
\begin{document}
\section*{UnsolvedMath ID: 30004815}
\textbf{Source problem code:} catalogue entry
\subsection{Definitions and mathematical statement}
Let \(X=X_1\cup X_2\cup X_3\) denote the supplied trisection of a smooth four-manifold, and put \(\Sigma=X_1\cap X_2\cap X_3\) and \(H_{ij}=X_i\cap X_j\) for distinct \(i,j\in\{1,2,3\}\). Here \(\Sigma\) is the central surface. A boundary-parallel arc tangle in a three-manifold is a finite collection of disjoint properly embedded arcs admitting mutually disjoint parallelism disks: each arc and an arc in the boundary cobound one such disk. A boundary-parallel disk tangle in a four-manifold is a finite collection of disjoint properly embedded disks admitting mutually disjoint parallelism three-balls: each disk and a disk in the boundary cobound one such three-ball, with corners allowed along their common boundary.

\textbf{Extracted bridge-position assertion.} The supplied record asserts that every smoothly embedded surface \(S\subset X\) can be isotoped to a surface \(S'\) such that, for some \(b\in\mathbb Z_{\geq 0}\),
\[
S'\pitchfork\Sigma,\qquad \#(S'\cap\Sigma)=2b,
\]
and, for every pair of distinct indices \(i,j\), the intersection \(S'\cap H_{ij}\) is a boundary-parallel arc tangle, while, for every \(i\in\{1,2,3\}\), the intersection \(S'\cap X_i\) is a boundary-parallel disk tangle. These conditions describe bridge position relative to the given trisection.

\textbf{Unresolved scope and intended target.} This is a transcription of the extracted assertion, not a recovered open conjecture. The supplied verification states that the report labels this assertion \emph{Theorem 5}, and that the neighboring open conjecture was not recovered. The record does not specify compactness, closedness, orientability, or boundary conditions for \(X\) and \(S\), nor whether a closed or relative trisection is intended. Consequently, the universal assertion above is schematic: the hypotheses needed for a precise theorem cannot be certified from the supplied extraction. No extension to surfaces with boundary, noncompact surfaces, or relative trisections is asserted here. The intended open problem remains unrecoverable.
\subsection{Short English statement}
The extracted assertion says that an embedded surface can be isotoped into bridge position for a given trisection: its intersections with the sectors are boundary-parallel disks, its intersections with pairwise overlaps are boundary-parallel arcs, and it meets the central surface transversely in an even number of points. The supplied verification identifies this assertion as Theorem~5, while the intended open conjecture and complete hypotheses are missing. The registry's open label therefore does not establish that this bridge-position assertion is an open problem.
\subsection{Sources}
\begin{itemize}
\item Local UnsolvedMath registry, \texttt{data/unsolvedmath/problems.json}: permanent ID \(30004815\), exact source problem code \texttt{catalogue entry}, title ``Bridge Position for Surfaces in Trisected Four-Manifolds.'' The supplied record belongs to the Oberwolfach Reports Open Problems collection. No individual proposer is supplied.
\item Supplied report citation: \emph{MATRIX-MFO Tandem Workshop: Invariants and Structures in Low-Dimensional Topology} (2022), Oberwolfach Reports, DOI \texttt{10.4171/owr/2021/42}, pp.~2327--2364. Supplied source URL and literature-source link, labeled ``OWR report'': \url{https://doi.org/10.4171/owr/2021/42}. This is the supplied authoritative report reference; its contents have not been independently consulted for this draft.
\item Supplied statement-verification note: the official report labels the bridge-position assertion Theorem~5; the extraction is a proved theorem rather than the neighboring open conjecture. This attribution is reproduced from the record, not independently verified.
\item Supplied literature assessment, checked exactly \texttt{2026-08-22}: status ``open.'' The assessment says that the extraction contains only the word ``conjecture'' and a theorem about bridge position, so the intended conjectural claim cannot be identified and remains open under the curation's triage rule. This is supplied dated curation, not an independently verified literature assessment. The background explicitly requests independent verification of time-sensitive status and cited claims.
\item The supplied proposed year and year tag are 2021, whereas the supplied report citation gives 2022. Both are retained without resolving their relationship.
\end{itemize}
\end{document}
